% =====================================================================
\chapter{Part A: System Model and Proposed Methods}
\label{ch:parta-methods}
\thispagestyle{plain}
% =====================================================================

\section{Introduction}

Part~A attacks the fronthaul load at its source and then makes the O-DU fast enough to work with the reduced data. The work falls into two threads. In the first (Sections~\ref{sec:a-model}--\ref{sec:rasbfp}) we compress the per-slot channel IQ using the structure it already has, namely a low numerical rank across antennas and sparsity in the delay domain. The second thread (Sections~\ref{sec:a-bench}--\ref{sec:ru-precoding}) is about fitting the per-slot precoder computation inside the 0.5~ms slot. It covers a screening benchmark of eight inversion methods including our learned inverters, the RMT-WS warm-start complex inversion method, and an analysis of where precoding should run to keep the fronthaul load small. The optimised C kernel that implements the precoder is described in Chapter~\ref{ch:impl}.

\section{System Model for IQ Compression}
\label{sec:a-model}

After OFDM demodulation, the O-RU holds for every slot the frequency-domain received matrix
\begin{equation}
  Y = H \odot X + W \in \mathbb{C}^{M\times N},
  \label{eq:a-model}
\end{equation}
where $M = 64$ is the number of antennas, $N = 1200$ the number of subcarriers (a 20~MHz carrier at numerology~0), $H$ the channel, $X$ the transmitted symbols and $W$ the noise. We draw the channel from the 3GPP TR 38.901 TDL-A delay profile~\cite{tr38901} and add exponential spatial correlation across antennas. If I and Q are 16-bit samples ($b = 16$), with $N_{\rm sym} = 14$ symbols per slot and $R_{\rm slot} = 1000$ slots/s, the raw antenna-domain fronthaul rate is
\begin{equation}
  R_{\rm raw} = 2\, M\, N\, N_{\rm sym}\, b\, R_{\rm slot} \approx 34.4~\text{Gbit/s per cell}.
\end{equation}
Every compression ratio in this report is measured against this 16-bit baseline.

\paragraph{Which IQ is compressed.} Delay-domain sparsity and low rank belong to the channel $H$, not to the data. When $X$ carries random QAM data, $H \odot X$ is neither sparse in delay nor low-rank, and so data-symbol IQ stays on O-RAN BFP. Our encoders target the IQ that carries channel information, i.e.\ the reference symbols (DMRS and SRS) and the channel matrices shipped for RU-side precoding (Section~\ref{sec:ru-precoding}). On a reference symbol the pilot sequence $X$ is known and has unit modulus. The O-RU can then de-rotate the received block, $Y \odot \bar X = H + \tilde W$, which is delay-sparse and low-rank. This channel-plus-noise matrix is what we compress in all experiments. Reference symbols are a small part of the fronthaul. An NR uplink slot of 14 OFDM symbols usually carries one or two DMRS symbols and at most one SRS symbol~\cite{ts38211}, so the encoders act on 1--3 of the 14 symbols and data symbols remain on BFP. Section~\ref{sec:ra-share} turns the compression ratios into a saving on the total fronthaul. The target NMSE of about $-20$~dB matches the noise floor at the 20~dB operating SNR. The compression error is thus of the same order as the estimation noise already in the channel estimate.

BFP ignores two kinds of structure in the signal:
\begin{itemize}
  \item \textbf{Low rank.} Because the antennas are correlated, the numerical rank $r$ of $Y$ is much smaller than $\min(M,N)$.
  \item \textbf{Delay-domain sparsity.} A TDL channel has just a few strong taps. After an inverse FFT along subcarriers, most of the energy of $Y$ therefore lies in a few delay bins.
\end{itemize}

\paragraph{Baseline: O-RAN BFP.} The O-RAN compression standard~\cite{oran_cus} groups the 12 subcarriers of a PRB and stores one shared exponent plus a $b$-bit mantissa for each I and Q value. The numbers become shorter, but there are just as many of them. The gain is about $16/b$, i.e.\ $\approx$2$\times$ at $b=8$.

Figure~\ref{fig:a-pipeline} shows where our encoders are placed. They sit in the O-RU after the FFT and replace (or wrap) BFP, while the decoder runs in the O-DU before equalisation.

\begin{figure}[H]
\centering
\begin{tikzpicture}[font=\footnotesize,>=Stealth,
   b/.style={draw,rounded corners,align=center,minimum height=1.0cm,inner sep=3pt},
   ours/.style={b,fill=good!15,draw=good!70!black,thick}]
  \node[b,fill=iitblue!12] (ant) {64 antennas\\RF + ADC};
  \node[b,fill=iitblue!12,right=0.45cm of ant] (fft) {CP removal\\+ FFT};
  \node[ours,right=0.45cm of fft] (enc) {\textbf{Encoder}\\CSEE / RAS-BFP};
  \node[b,fill=gray!10,right=2.4cm of enc] (dec) {Decoder\\$\hat Y$};
  \node[b,fill=good!8,right=0.45cm of dec] (eq) {Channel est.,\\equalisation,\\decoding};
  \draw[->,thick] (ant) -- (fft);
  \draw[->,thick] (fft) -- node[above]{$Y$} (enc);
  \draw[->,line width=2pt,accent] (enc) -- node[above,text=accent]{fronthaul} node[below,text=accent,font=\scriptsize,align=center]{4--31$\times$ fewer\\reference-symbol bits} (dec);
  \draw[->,thick] (dec) -- (eq);
  \node[draw,dashed,rounded corners,fit=(ant)(fft)(enc),inner sep=6pt,label={[font=\footnotesize\bfseries]above:O-RU}] {};
  \node[draw,dashed,rounded corners,fit=(dec)(eq),inner sep=6pt,label={[font=\footnotesize\bfseries]above:O-DU}] {};
\end{tikzpicture}
\caption{Placement of the proposed structure-aware IQ encoders in the 7-2x uplink chain.}
\label{fig:a-pipeline}
\end{figure}

\section{CSEE: Delay-Domain Sparse Encoder}
\label{sec:csee}

CSEE moves each antenna row of $Y$ into the delay domain. It then keeps the strongest taps and quantises them, in four steps.
\begin{enumerate}
  \item $Y_d = \mathrm{IFFT}_N(Y)$ (row-wise, along subcarriers).
  \item Select the support $S = \mathrm{TopK}(|Y_d|)$ of the $K$ largest-magnitude delay coefficients.
  \item Quantise the kept coefficients with per-antenna BFP of step $\Delta$ and transmit values and support.
  \item At the O-DU, reconstruct $\hat Y = \mathrm{FFT}_N(Y_d|_S)$.
\end{enumerate}

\begin{theorem}[CSEE error bound]
Let $\sigma^2_{\rm tail}$ be the energy of the discarded delay coefficients and $\sigma_Y^2$ the energy of $Y$. With a uniform quantiser of step $\Delta$ on the $K$ kept coefficients,
\begin{equation}
  \mathrm{NMSE}_{\rm CSEE} \le \frac{\sigma^2_{\rm tail}}{\sigma^2_Y} + \frac{\Delta^2}{12}\,\frac{K}{MN\sigma_Y^2}.
\end{equation}
\end{theorem}
The first term is the truncation error and the second is the quantisation error. We control both through $K$ and the mantissa width. Since the FFT is unitary, an error in the delay domain is the same as the error in the frequency domain. The FFTs dominate the cost, which is $O(MN\log N)$.

\section{RAS-BFP: Randomised Sketch and QR Encoder}
\label{sec:rasbfp}

RAS-BFP builds a low-rank factorisation of $Y$ from a randomised sketch. This costs far less than an SVD. The steps are as follows.
\begin{enumerate}
  \item Sketch with a subsampled randomised Hadamard transform (SRHT): $Z = \Omega Y \in \mathbb{C}^{r\times N}$.
  \item Thin QR of the sketch: $Z^H = QR$, with $Q \in \mathbb{C}^{N\times r}$ orthonormal.
  \item Project: $L = YQ \in \mathbb{C}^{M\times r}$.
  \item Transmit $\mathrm{BFP}(L)$ and $\mathrm{BFP}(Q)$; the O-DU reconstructs $Y \approx LQ^H$.
\end{enumerate}

\begin{theorem}[RAS-BFP approximation, after Halko--Martinsson--Tropp~\cite{halko2011}]
With an SRHT sketch of $r$ rows (plus oversampling) and $Y_r$ the best rank-$r$ approximation of $Y$,
\begin{equation}
  \|Y - LQ^H\|_F^2 \le (1+\varepsilon_r)\,\|Y - Y_r\|_F^2 \quad\text{with probability at least } 1-\delta.
\end{equation}
\end{theorem}

\paragraph{Rank under a bit budget.} Together, $L$ and $Q$ hold $r(M+N)$ complex values. With $b$-bit mantissas for I and Q the payload comes to $2b\,r(M+N)$ bits, plus one exponent per block. Given a total budget $B_{\rm total}$, the largest rank we can afford is
\begin{equation}
  r_{\max} = \left\lfloor \frac{B_{\rm total}}{2b\,(M+N)} \right\rfloor .
\end{equation}
The truncation error in the theorem above decreases with $r$. The quantisation error decreases as $2^{-2b}$. On the budget line $r\,b = B_{\rm total}/(2(M+N))$ the encoder picks the pair $(r, b)$ with the smallest sum of the two terms. It is a one-dimensional search over $r \le r_{\max}$ and is done once per configuration.

RAS-BFP costs $O(rMN)$, compared with $O(MN\min(M,N))$ for a truncated SVD, and it achieves the same ratio at the $-20$~dB NMSE target. At link level the picture is different. There it needs twice the rank of SVD before it becomes lossless, which puts its operating point at 4$\times$ (Table~\ref{tab:a-methods} and Section~\ref{sec:ra-link}).

\begin{table}[H]
\caption{IQ compression methods compared in Part~A.}
\label{tab:a-methods}
\centering
\resizebox{\textwidth}{!}{%
\begin{tabular}{lllll}
\toprule
Method & Structure exploited & Complexity & Theory & \makecell{Link-level\\operating point} \\
\midrule
O-RAN BFP & dynamic range & $O(MN)$ & empirical quantisation & 2$\times$ \\
Truncated SVD & low rank & $O(MN\min(M,N))$ & Eckart--Young~\cite{eckart1936} & 8$\times$ \\
\ours \textbf{CSEE (ours)} & delay sparsity & $O(MN\log N)$ & Theorem~3.1 & \textbf{31$\times$} \\
\ours \textbf{RAS-BFP (ours)} & low rank & $O(rMN)$ & Theorem~3.2 & \textbf{4$\times$} \\
\bottomrule
\end{tabular}}
\end{table}

\section{Matrix-Inversion Benchmark and Learned Inverters}
\label{sec:a-bench}

MMSE equalisation, ZF/RZF precoding, channel estimation and our own encoders each need a matrix inverse or a decomposition \emph{every slot}. Before settling on a family of methods for a 0.5~ms budget, we ran a screening benchmark of eight methods on the CPU.

\paragraph{Dataset.} We use matrices $A = G/\sqrt{n} + 2I$ with $G_{ij}\sim\mathcal N(0,1)$ and $n \in \{10, 100, 500\}$, with 1000 matrices per size in float32. These matrices are real and well conditioned, and they were kept simple on purpose. The screening only compares method \emph{families} under identical conditions. The selected family is later tested on the real complex RZF workload (Section~\ref{sec:rzf}) and as a C kernel (Chapter~\ref{ch:impl}). The learned models are trained on 80\,\% of the matrices, and timing is measured on 200 held-out matrices, one at a time after warm-up, on the CPU only. We report the mean, median and 95th-percentile time.

\paragraph{Classical methods (6).} These are LAPACK \texttt{getri} called through numpy, LU (scipy), QR, SVD, Gauss--Jordan elimination and the Newton--Schulz iteration.

\paragraph{Our learned inverters (2).}
\begin{itemize}
  \item \textbf{InverseNet-NS}: the Newton--Schulz iteration unrolled for 8 steps,
  \[ X_{k+1} = X_k(\beta_k I - \gamma_k A X_k), \qquad X_0 = \alpha A^T/(\|A\|_1\|A\|_\infty), \]
  where the step coefficients $(\beta_k,\gamma_k)$ and the scale $\alpha$ of the initial guess are learned.
  \item \textbf{InverseNet-Ultra}: a small hyper-network predicts the initial guess $X_0 = a_s A^T/\|A\|_F^2 + i_s I$ from four matrix statistics (trace, Frobenius, 1- and $\infty$-norm). It is followed by 6 learned Newton--Schulz steps. Thanks to the better initial guess, two steps can be dropped.
\end{itemize}
Both learned iterations keep the mathematical structure of a provably convergent method and learn only its coefficients. As a result they need very few parameters and generalise across matrices. Their cost is a fixed number of batched matrix multiplications, and GPUs, NPUs and tensor cores run that operation faster than any other.

\section{The Real Workload: Complex RZF Precoding}
\label{sec:rzf}

In the downlink the O-DU (or the O-RU, Section~\ref{sec:ru-precoding}) computes for every PRB and every slot the regularised zero-forcing precoder
\begin{equation}
  G = HH^H + \alpha I \in \mathbb{C}^{K\times K}, \qquad W = H^H G^{-1} \in \mathbb{C}^{M\times K},
  \label{eq:rzf}
\end{equation}
for $K$ users and $M$ antennas, 273 times inside a 0.5~ms slot. To test inversion methods on this workload we wrote a realistic complex channel simulator (\texttt{cmatinv/}). It models 3GPP TR 38.901 UMa NLOS path loss and shadowing, which gives a large near-far spread of user powers. It adds the TDL-A 23-tap profile over 100~MHz, per-user angular-spread spatial correlation, an optional Rician line-of-sight component, and Clarke/Jakes channel aging $\rho = J_0(2\pi f_D T)$ for user speeds of 3--120~km/h. We simulate $M \in \{64, 128\}$ and $K \in \{8, 16, 32\}$ with 273 PRBs, 30~kHz subcarrier spacing and a 3.5~GHz carrier.

\section{RMT-WS: Random-Matrix-Guided Warm-Start Inversion}
\label{sec:rmtws}

RMT-WS inverts the complex Hermitian Gram matrices of Eq.~\eqref{eq:rzf} using matrix multiplications alone. Its starting point is a near-optimal initial guess obtained from random-matrix theory. When the channel has changed only slightly, it re-uses the inverse from the previous slot instead. Figure~\ref{fig:rmtws} shows the flow for one PRB.

\begin{figure}[htbp]
\centering
\begin{tikzpicture}[font=\footnotesize,>=Stealth,
   s/.style={draw,rounded corners,fill=soft,align=center,text width=3.4cm,minimum height=0.95cm},
   d/.style={draw,diamond,aspect=2.2,fill=orange!15,align=center,inner sep=1pt,font=\scriptsize},
   g/.style={s,fill=good!15}]
  \node[s] (g)    at (0,0)     {Gram $G = HH^H + \alpha I$};
  \node[s] (eq)   at (0,-1.4)  {Equilibrate\\$\hat G = D^{-1/2} G D^{-1/2}$};
  \node[d] (dec)  at (0,-3.0)  {previous $X$\\available?};
  \node[s] (int)  at (-3.6,-4.6) {Spectral interval $[a,b]$\\moment-matched MP\\($K^2$ work)};
  \node[s] (cheb) at (-3.6,-6.4) {Chebyshev polynomial\\on $[a,b]$ (cold start)};
  \node[s] (warm) at (3.6,-4.6)  {Warm start: $X_0 = X_{t-1}$\\rescaled by new $D$};
  \node[d] (cert) at (3.6,-6.4)  {certificate\\$\hat r \le 0.3$?};
  \node[g] (ns)   at (0,-8.4)  {Newton--Schulz\\$X \leftarrow X(2I - \hat G X)$\\steps from $\hat r$};
  \node[s] (out)  at (0,-10.1) {$G^{-1} = D^{-1/2} X D^{-1/2}$\\$\Rightarrow W = H^H G^{-1}$};
  \draw[->,thick] (g) -- (eq);
  \draw[->,thick] (eq) -- (dec);
  \draw[->,thick] (dec.west) -| node[pos=0.25,above]{no} (int.north);
  \draw[->,thick] (dec.east) -| node[pos=0.25,above]{yes} (warm.north);
  \draw[->,thick] (int) -- (cheb);
  \draw[->,thick] (warm) -- (cert);
  \draw[->,thick] (cert.west) -- node[above]{no: cold start} (cheb.east);
  \draw[->,thick] (cheb.south) |- (ns.west);
  \draw[->,thick] (cert.south) |- node[pos=0.25,right]{yes} (ns.east);
  \draw[->,thick] (ns) -- (out);
\end{tikzpicture}
\caption{RMT-WS inversion of one PRB: equilibration, random-matrix spectral interval, Chebyshev cold start or certified warm start, and Newton--Schulz refinement.}
\label{fig:rmtws}
\end{figure}

\paragraph{1. Equilibration.} Setting $\hat G = D^{-1/2} G D^{-1/2}$ with $D = \mathrm{diag}(G)$ gives a unit diagonal. It also removes the near-far power spread, which would otherwise ruin every series-based method.

\paragraph{2. Eigensolve-free spectral interval.} For an i.i.d.\ channel the eigenvalues of $\hat G$ follow the Marchenko--Pastur (MP) law~\cite{marchenko1967} (Figure~\ref{fig:mp} in Chapter~\ref{ch:parta-results}). Correlation and line-of-sight make the spectrum wider, and a fixed-ratio MP edge $(1\mp\sqrt{K/M})^2$ cannot follow this. We match the first two spectral moments instead. The spectral mean of $\hat G$ is 1 and its variance is $v = \|\hat G\|_F^2/K - 1$, from which
\begin{equation}
  [a,b] = 1 + v/\bar s \mp 2\sqrt{v},
\end{equation}
follows. Computing it takes only $K^2$ operations and no eigen-decomposition. Correlated and LoS channels raise $v$, so the interval widens by itself.

\paragraph{3. Cold start.} A Chebyshev polynomial of degree $d$ on $[a,b]$ provides an initial inverse $X_0 = p_d(\hat G)$. Its residual $\|I - \hat G X_0\|$ is minimax-optimal on the interval. From there, Newton--Schulz steps $X \leftarrow X(2I - \hat G X)$ converge quadratically and need nothing but GEMMs.

\paragraph{4. Warm-start tracking.} At pedestrian speed the channels of consecutive slots are strongly correlated. The previous inverse, rescaled by the new diagonal, is then a very good starting point and needs just one or two Newton steps.

\paragraph{5. Residual certificate.} We use two Gaussian probe vectors. Evaluating $\hat r = \|(I - \hat G X)z\|/\|z\|$ then costs $4K^2$ multiply-accumulates, or about $4/K$ of a GEMM. The estimate predicts how many Newton steps are needed, $\lceil \log_2(\ln\varepsilon / \ln \hat r)\rceil$. If $\hat r > 0.3$, the warm start is declared non-contractive and the cold start runs in its place. With this check RMT-WS is safe at every user speed.

\paragraph{Doppler-to-warm-start law.} Expanding the AR(1) aging channel to first order gives the initial residual of the warm start as
\begin{equation}
  r_0 \approx c\cdot 2\sqrt{\frac{K}{M-K}}\,\sqrt{1-\rho}, \qquad 1-\rho \approx (\pi f_D T)^2,
  \label{eq:doppler}
\end{equation}
and $r_0$ therefore grows linearly with the user speed. We use Eq.~\eqref{eq:doppler} as a design rule. For a given array and user count, the O-DU knows in advance up to which speed the warm start pays off (about 39~km/h for $M=64$, $K=16$; 60~km/h for $K=8$; 23~km/h for $K=32$).

\paragraph{Certificate guarantee.} With $R = I - \hat G X$, the probe estimate $\hat r^2$ is an unbiased estimate of $\|R\|_F^2/K$. Applying a Laurent--Massart tail bound for weighted $\chi^2$ sums, we get
\begin{equation}
  \Pr\big(\hat r^2 \le (1-d)\, r^2\big) \le e^{-p d^2 s/2},
\end{equation}
where $p$ is the number of probes and $s$ is the stable rank of $R$. We audited 480 PRB cases between 3 and 120~km/h and found no accepted warm start that was non-contractive.

\paragraph{Why RMT-WS suits the O-DU.} The method adjusts by itself to correlated and near-far channels and needs no eigen-decomposition. It is built from matrix multiplications only, which fits GPU, NPU and tensor-core hardware well. Its cost scales with the user speed, and it comes with its own accuracy certificate.

\section{Where Should Precoding Run? RU-Side Precoding}
\label{sec:ru-precoding}

Where the precoder runs has a large effect on the fronthaul rate (Figure~\ref{fig:ru-prec}). When the O-DU precodes, all $M=64$ antenna streams of IQ samples have to cross the fronthaul. When the O-RU precodes instead (the O-RAN ``category B'' radio), the fronthaul only carries the $K=16$ layer streams and the beamforming weights $W$. Under mobility the weights must be refreshed every slot, so the architecture only works with a fast per-slot precoder. RMT-WS and the fused kernel of Chapter~\ref{ch:impl} provide one.

\begin{figure}[H]
\centering
\begin{tikzpicture}[font=\footnotesize,>=Stealth,
   b/.style={draw,rounded corners,align=center,minimum height=0.9cm,inner sep=3pt,minimum width=2.9cm}]
  \node[b,fill=good!15] (du1) at (0,0) {O-DU\\layers + \textbf{precoding}};
  \node[b,fill=iitblue!15] (ru1) at (8.6,0) {O-RU\\64 antennas};
  \draw[->,line width=4pt,accent] (du1) -- node[above,text=accent]{$M = 64$ antenna streams} node[below,text=accent]{118.4 Gbit/s} (ru1);
  \node[anchor=east,font=\footnotesize\bfseries] at (-1.6,0) {(a)};
  \node[b,fill=good!15] (du2) at (0,-2.0) {O-DU\\layers + weights $W$};
  \node[b,fill=iitblue!15] (ru2) at (8.6,-2.0) {O-RU\\\textbf{precoding} + 64 ant.};
  \draw[->,line width=1.6pt,good!70!black] (du2) -- node[above,text=good!50!black]{$K = 16$ layers + 8-bit weights / slot} node[below,text=good!50!black]{32.8 Gbit/s ($-72\,\%$)} (ru2);
  \node[anchor=east,font=\footnotesize\bfseries] at (-1.6,-2.0) {(b)};
\end{tikzpicture}
\caption{(a) DU-side precoding sends every antenna stream over the fronthaul; (b) RU-side precoding sends only the layer streams plus per-slot weights, cutting the downlink fronthaul rate by 72\,\% (100~MHz, $M=64$, $K=16$, BFP-9).}
\label{fig:ru-prec}
\end{figure}

\section{Summary}

This chapter set out the system model and methods of Part~A. We introduced two structure-aware encoders with error bounds, CSEE and RAS-BFP, and an eight-method inversion screening benchmark with learned inverters. We then described the RMT-WS warm-start complex inversion with its Doppler law and certificate, and the RU-side precoding architecture that fast per-slot precoding makes possible. The simulation results for these methods are given in Chapter~\ref{ch:parta-results}.
